%-------------------------------------------------------------------------------
\section{Security Properties} \label{sec:security_properties}
%-------------------------------------------------------------------------------

In this section, we formally define the security properties of third-party online payment protocols.
The core security goal of protocols is to ensure that the user, the merchant, and the payment service provider reach a consistent view of a transaction.
The protocol must guarantee that (a) all parties agree on the payment order being processed, (b) each payment is authorized by the user, and (c) the reported payment result faithfully reflects the actual fund transfer.

All properties considered in this section are instances of \emph{injective agreement}~\cite{Lowe1997AuthenticationHierarchy}.
Injective agreement strengthens standard agreement by requiring a one-to-one correspondence between protocol runs: 
each successful local completion must correspond to a unique execution of the peer, ruling out replay, duplication, or reuse of protocol messages.
Two types of atomic facts are necessary to define injective agreement:
\begin{enumerate}[leftmargin=*]
  \item $\mathsf{Running}(A,B,x)$: 
    records that party $A$ is interacting with $B$ in a concrete execution of the protocol and that the interaction involves data $x$ (an objective execution fact);
  \item $\mathsf{Commit}(A,B,x)$:
    denotes that party $A$ subjectively believes that it has reached agreement with $B$ on data $x$ (e.g., after receiving and accepting a message).
\end{enumerate}
This notion is particularly well-suited for online payment scenarios, where both consistency and uniqueness of transactions are fundamentally essential requirements.
We structure the security properties according to the logical progression of a payment protocol, namely agreement on the \emph{order}, agreement on \emph{user authorization}, and agreement on the \emph{result}.

Each property instantiates the following template, in which $\mathsf{Commit}$ and $\mathsf{Running}$ are replaced by the fact pair of that property:
\begin{tightflalign}
  \TAll \; \text{\Tterm{a}} & \text{\Tterm{ b x \#i.}} \; \text{\TFact{Commit(a, b, x)@\#i}} &\\
  \implies & \TEx \; \text{\TLemma{\#j.}} \; \text{\TFact{Running(b, a, x)@\#j}} \; \land\\
  \lnot \; &(\;\TEx \; \text{\TLemma{a2 b2 \#i2.}} \; \text{\Tterm{\#i2}} \neq \text{\Tterm{\#i}} \; \land\\
  & \hspace{2em} \text{\TFact{Commit(a2, b2, x)@\#i2}})
\end{tightflalign}
The first conjunct requires that the peer actually ran the protocol on the same data, and the second requires that no other commit uses that data, which rules out replay and duplication.
The Commit and Running facts follow the standard interpretation introduced in~\cite{Lowe1997AuthenticationHierarchy}.

\textbf{Injective Agreement on Order.}
This property ensures that all involved parties process the same payment order and that each order is handled at most once; a violation indicates that an adversary has tampered with the order information.
Since it involves three principals, we consider \emph{PS} and \emph{MS}, and \emph{U} and \emph{MS}, separately.
\emph{IA-O1~(\emph{PS} $\leftarrow$ \emph{MS})} uses \TFact{PSCommit(ps, ms, order)} and \TFact{MSRunning(ms, ps, order)}: if \emph{PS} accepts a payment request for an order allegedly issued by \emph{MS}, then \emph{MS} has initiated a request for that order, and no other request for the same order is accepted.
The term \Tterm{order} includes the merchant's order identifier, the payment amount, and the product description.
\emph{IA-O2~(\emph{U} $\leftarrow$ \emph{MS})} uses \TFact{UMCommit(u, ms, goods)} and \TFact{MSRunning(ms, u, goods)}: the user and the merchant agree on the purchased goods, and each order is processed exactly once.
IA-O2 requires agreement on the goods only, since the user is primarily concerned with what is being bought.

\textbf{Injective Agreement on User Authorization.}
This property ensures that each successful payment is backed by a unique and explicit authorization from the user, which prevents both unauthorized payments and the reuse of one authorization for multiple transactions.
\emph{IA-UA~(\emph{PS} $\leftarrow$ \emph{U})} uses \TFact{PaymentComplete(u, h(p), s)} and \TFact{UserInputPwd(u, p, s)} with the session \Tterm{s} as the agreed data: if \emph{PS} completes a payment based on the hash \Tterm{h(p)} of a password, then \emph{U} has input \Tterm{p} in session \Tterm{s}, and each password input authorizes at most one payment.

\textbf{Injective Agreement on Result.}
This property guarantees that successful payment notifications cannot be forged, replayed, or duplicated, so that an adversary cannot convince the user or the merchant that a payment has been completed when no fund transfer has occurred.
\emph{IA-R1~(\emph{MS} $\leftarrow$ \emph{PS})} uses \TFact{MSCommit(ms, ps, order)} and \TFact{PSRunning(ps, ms, order)}: a successful notification received by \emph{MS} has been generated by \emph{PS}, and \emph{MS} receives at most one such notification per order.
\emph{IA-R2~(\emph{U} $\leftarrow$ \emph{PS})} uses \TFact{UPCommit(u, ps, goods)} and \TFact{PSRunning(ps, u, goods)}: a successful notification observed by \emph{U} has been sent by \emph{PS}, and each order produces at most one notification.
